% !TeX program = xelatex
% UTF-8. Compile from this directory or from the repository root with XeLaTeX.
% Source: Google Drive / Notebook / respiratory diseases of neonates.
% pic01.jpg = 20261016_110645924.jpg; Drive ID 1mzyPKep99yx2P8iyKedhZINc7_HtEnf8.
% pic02.jpg = 20261016_110651737.jpg; Drive ID 1rKQgj9HRKk6L4BJxEh_Fj9XDHlqhlHzD.
% pic03.jpg = 20261016_110656297.jpg; Drive ID 1LCVh8KNRviZs_2p-PZGoBPRQI-xb3Zho.
% pic04.jpg = 20261016_110702472.jpg; Drive ID 1P4EL-Xl3bB5pkfCsTym74rRDzCP3r17m.
% pic05.jpg = 20261016_110706259.jpg; Drive ID 1XkJDnw97ExR-tPjYf7-3xYXGQ4XOnfTk.
% pic06.jpg = 20261016_110711935.jpg; Drive ID 1yD3hNZXvd1_0DyIvpzZ0CCuTcPEy9iRI.
% pic07.jpg = 20261016_110716058.jpg; Drive ID 1Rnlmc6HSVirGKMk81Uv0vCRUkbyzRMRM.
% pic08.jpg = 20261016_110721100.jpg; Drive ID 1fz_EROsgvV7DAD64AgxksEp2Ft_PXT3_.
% pic09.jpg = 20261016_110725492.jpg; Drive ID 1tSMFumvh4BREgs61fHGXJV0lMH9xKF7r.
% pic10.jpg = 20261016_110732121.jpg; Drive ID 1SCIvFqmlA7vcQ8Ao6mnbpQQwuS7oPB35.
% pic11.jpg = 20261016_110736293.jpg; Drive ID 1unZIcXQMp61sVSpvs7GvTIB5_hOJTb8M.
% pic12.jpg = 20261016_110742525.jpg; Drive ID 1aL0huyTqZymcwVzBUTKls0P7W4pemuFj.
% pic13.jpg = 20261016_110746796.jpg; Drive ID 1AOLtYowXgbM0t3FtfiPr98dgE5DLmsc9.
% pic14.jpg = 20261016_110752403.jpg; Drive ID 17hkdsLwNY_b0ZJ5APWgFI2F8tzlmELWa.
% Images are numbered in ascending order of their original filenames.
% Define \NeonatalRespiratoryInNotebook before input into an existing XePersian document.
\ifdefined\NeonatalRespiratoryInNotebook
  \let\NRfinish\relax
\else
  \def\NRfinish{\end{document}}
  \documentclass[12pt,a4paper]{article}
  \usepackage[margin=22mm]{geometry}
  \usepackage{amsmath,amssymb}
  \usepackage{graphicx}
  \usepackage{enumitem}
  \usepackage{xepersian}
  \IfFontExistsTF{Amiri}{\settextfont{Amiri}}{\settextfont{DejaVu Sans}}
  \setlatintextfont{DejaVu Serif}
  \setlength{\parindent}{0pt}
  \setlength{\parskip}{5pt}
  \setlength{\emergencystretch}{3em}
  \begin{document}
\fi
\providecommand{\NRmark}[1]{\textbf{#1}}
\providecommand{\NRunclear}[1]{\textbf{[خوانش نامطمئن: #1]}}
\providecommand{\NRmath}[1]{\lr{\(#1\)}}
\providecommand{\NRpage}[1]{\clearpage\section*{تصویر #1}}

\begin{center}
{\Large\bfseries ۶۱ ـ بیماری‌های تنفسی نوزادان}\par
\lr{Respiratory diseases of neonates}
\end{center}
\textbf{یادداشت رونویسی:}
متن، اعداد، عبارات انگلیسی و شاخه‌های دست‌نویس به ترتیب تصاویر بازنویسی شده‌اند.
عبارت‌های هایلایت‌شده پررنگ‌اند؛ خوانش نامطمئن صریحاً مشخص شده است.
تصاویر کامل در پیوست همین سند آمده‌اند.
سربرگ چاپی همهٔ برگه‌ها: \lr{Subject:}، \lr{Year:}، \lr{Month:} و \lr{Date:}؛
خانه‌های تاریخ خالی‌اند. عنوان دست‌نویس تصاویر ۱ تا ۳ «۶۱ ـ بیماری‌های تنفسی نوزادان»
و سربرگ تصاویر ۴ تا ۱۴ «۶۱» است.
شماره‌های چاپی کنار سطرها ۱ تا ۹ و شماره‌های دست‌نویس پایین صفحات ۱ تا ۱۴ هستند.

\NRpage{۱}
\subsection*{دیسترس تنفسی نوزاد}
\begin{itemize}
\item علائم دیسترس تنفسی نوزاد: تاکی‌پنه، رتراکشن بین‌دنده‌ها، کاهش تبادل هوا، سیانوز، گرانتینگ بازدمی، نازال فلرینگ.
\item دیسترس تنفسی در اختلالات: ریوی، قلبی، هماتولوژیک، عفونی، آناتومیک، متابولیک.
\item \lr{RDS}: به علت سورفاکتانت دفیشنسی؛ \NRmark{سیانوز + تاکی‌پنه}.
\item پنومونی: درگیری لوبار یا اینترستیشیال.
\item مکونیوم آسپیریشن: پنومونیت شیمیایی ایجاد می‌کند + \NRmark{هایپوکسی} + \NRmark{فشار خون ریوی}.
\item هیدروپس فتالیس: \NRmark{آنمی + هایپوآلبومینمی} ایجاد می‌کند؛ \lr{HF} با برون‌ده قلبی بالا + ادم ریوی.
\item پولمونری هایپوپلازی: دو نوع همراه با تولد و اکتسابی؛ \NRmark{فشار خون ریوی}.
\item در همهٔ بیماران دیسترس تنفسی \NRmark{\lr{ABG}} باید مانیتور شود؛ به دنبال توجه به وجود یا عدم سیانوز.
\end{itemize}
\subsection*{\lr{ABG} نرمال در نوزادان}
\begin{itemize}
\item در نوزادان زودرس عدد پایین‌تر است.
\item \NRmath{PaO_2}: دقیقهٔ ۳: \lr{55--60 mmHg}؛ ساعت ۴: \lr{70 mmHg}؛ ساعت ۲۴: \lr{90 mmHg}.
\item \NRmath{PaCO_2}: \lr{35--40 mmHg}.
\item \lr{pH}: \lr{7.35--7.40}.
\end{itemize}
\subsection*{چه زمانی سیانوز قابل مشاهده است؟}
\begin{itemize}
\item \NRmark{\lr{5 gr}} هموگلوبین اشباع‌نشده وجود داشته باشد.
\item اگر آنمی وجود داشته باشد ممکن است سیانوز را نبینیم.
\item همچنین زردی.
\end{itemize}

\NRpage{۲}
\subsection*{اسیدوز متابولیک در نوزادان}
\begin{itemize}
\item \NRmath{pH<7.25}، \NRmath{bicarb<18\,\mathrm{mEq/L}} و \NRmath{PaCO_2} نرمال یا کاهش‌یافته.
\item اپروچ اسیدوز متابولیک: علت را مشخص کرد؛ درمان پاتوفیزیولوژی.
\item دوز بی‌کربنات در صورت نیاز: \lr{1 mEq/kg/dose, slow IV}.
\item قبل از تجویز باید \NRmath{PaCO_2} نزدیک نرمال یا پایین باشد.
\item \NRmark{در اسیدوز تنفسی بی‌کربنات نمی‌دهیم.}
\item اگر مادهٔ نارکوتیک از طریق جفت وارد بدن جنین شده و نوزاد دچار ساپرشن تنفسی شده:
\NRmark{اول ونتیلاسیون برقرار شود، سپس نالوکسان وارد شود.}
\end{itemize}
\subsection*{اتیولوژی‌های شایع دیسترس تنفسی نوزاد}
\begin{itemize}
\item نوزاد پره‌ترم: \NRmark{\lr{RDS*}}، اریتروبلاستوزیس فتالیس، هیدروپس غیرایمیون، خونریزی مغزی.
\item نوزاد ترم: پنومونی\lr{*}، مکونیوم آسپیریشن، پلی‌سیتمی، آسپیریشن مایع آمنیوتیک؛
\lr{Persistent pulmonary hypertension of the neonates*}.
\item نوزاد پره‌ترم و ترم: \NRmark{سپسیس \lr{GBS*}}، \NRmark{\lr{TTN*}}، پنوموتوراکس خودبه‌خودی، آنومالی‌ها، بیماری قلبی مادرزادی، هایپوپلازی پولمونری، نقص‌های متابولیک، عفونت ویروسی.
\item آنومالی‌ها: آمفیزم لوبار.
\end{itemize}
\subsection*{اقدامات تشخیصی برای نوزاد با دیسترس}
\NRmark{\lr{CXR, ABG, CBC, BC, BG, ECG}}.
\begin{itemize}
\item \lr{ECG}: در صورت سوفل، کاردیومگالی، هایپوکسی راجعه.
\item شک به اختلالات ساختاری قلب، \lr{PPHN}.
\end{itemize}

\NRpage{۳}
\begin{itemize}
\item \lr{BG}: \NRmark{هایپوگلایسمی می‌تواند دیسترس تنفسی بدهد.}
\item هایپرگلایسمی می‌تواند در همراهی با و یا علت دیسترس تنفسی باشد \NRunclear{واژهٔ نخست؛ احتمال خوانش «هایپوگلایسمی» نیز وجود دارد}.
\item \lr{CXR}: پنومونی، کاردیومگالی، پنوموتوراکس، آنومالی‌ها.
\end{itemize}
\subsection*{\lr{RDS} (هایالن ممبرین دیزیز)}
\begin{itemize}
\item سورفاکتانت کافی وجود ندارد.
\item دو نوع سلول: ۹۰\% سلول‌های تیپ \lr{I} و ۱۰\% سلول‌های تیپ \lr{II}.
\item تیپ \lr{II} از هفتهٔ ۲۰: سورفاکتانت.
\item سورفاکتانت از \NRmark{آتلکتازی} جلوگیری می‌کند.
\item در سن ۳۴–۳۶ هفتگی سطح سورفاکتانت به حداقل مورد نیاز می‌رسد.
\item سورفاکتانت وارد مایع آمنیوتیک می‌شود.
\item سطح \NRmark{لسیتین} مایع آمنیوتیک بلوغ ریهٔ جنین را نشان می‌دهد.
\item \NRmark{\NRmath{L/S\geq 2/1}}: \NRmark{بلوغ ریه}؛ \lr{S}: اسفنگومیلین.
\item اگر \lr{L/S} بوردرلاین بود: سطح \NRmark{فسفاتیدیل گلیسرول} را می‌سنجیم.
\item \NRmark{مادر \lr{DM} بود}: بلوغ ریه را تأخیر می‌اندازد.
\end{itemize}
\subsection*{علائم کمبود سورفاکتانت}
آتلکتازی، کاهش \lr{FRC}، هایپوکسمی شریانی، دیسترس تنفسی.
\subsection*{سایر علل کاهش تولید سورفاکتانت}
هایپوولمی، هایپوترمی، اسیدوز، هایپوکسی، اختلالات ژنتیکی نادر.
\par کمبود سورفاکتانت منجر به شانت پولمونری می‌شود.

\NRpage{۴}
\begin{itemize}
\item آتلکتازی: دشوار شدن تدریجی اکسپنشن ریه؛ پس نوزاد توانایی زیادی ندارد.
\item نوزاد فشار دم خود را بالا می‌برد؛ رتراکشن ایجاد می‌شود.
\item \NRmark{نوزادی که رتراکشن دارد، با احتمال بالا آتلکتازی هم دارد.}
\item کاهش کمپلیانس ریه + افزایش فشار دم (تاکی‌پنه): تبادل ضعیف هوا، افزایش فضای مردهٔ فیزیولوژیک، آلوئولار هایپوونتیلیشن، هایپرکاپنه.
\item هایپوکسی، هایپرکاپنه + اسیدوز: کاهش تولید سورفاکتانت؛ چرخهٔ معیوب؛ پولمونری هایپرتنشن.
\item کیس تیپیک \lr{RDS}: \NRmark{نوزاد نارس با \lr{L/S} پایین}.
\end{itemize}
\subsection*{سایر ریسک فاکتورهای \lr{RDS}}
سابقهٔ قبلی نوزاد نارس یا \lr{RDS}، دیابت مادر، هایپوکسی، دیسترس جنینی، آسفیکسی، پسر بودن، قل دوم، بارداری \lr{twin}، زایمان با \lr{C/S} بدون لیبر.
\begin{itemize}
\item \lr{RDS} \NRmark{بلافاصله} در اتاق زایمان رخ می‌دهد؛ نوزادان بسیار زودرس \NRmark{۲۸–۳۰ \lr{w}}.
\item \NRmark{پس از ۳–۴ ساعت} از تولد: نوزاد \lr{34 w}؛ ذخیرهٔ اولیه از سورفاکتانت وجود دارد.
\end{itemize}
\subsection*{تظاهرات بالینی \lr{RDS}}
\NRmark{سیانوز، تاکی‌پنه، نازال فلرینگ، رتراکشن جناغی و بین‌دنده‌ای، گرانتینگ}.
\par گرانتینگ ناشی از بسته شدن گلوت طی بازدم.
\subsection*{\lr{CXR}}
\begin{itemize}
\item آتلکتازی: \lr{ground glass haze in the lung surrounding air filled bronchi (}\NRmark{\lr{Air bronchogram}}\lr{)}.
\item آتلکتازی خیلی شدید: \lr{Airless lung field (}\NRmark{\lr{whiteout}}\lr{)}.
\end{itemize}

\NRpage{۵}
\begin{itemize}
\item در صورت عدم درمان \lr{RDS}، طی ۷۲ ساعت اول: افزایش دیسترس و هایپوکسمی.
\item اندیکاسیون ونتیلاسیون کمکی: \lr{RDS} شدید که \NRmark{ادم}، آپنه و \NRmark{نارسایی تنفسی} ایجاد کند.
\item \lr{RDS} غیرعارضه‌دار خودبه‌خودی بهبود می‌یابد؛ شواهد بهبودی: \NRmark{دیورز، دفع سریع ادم}.
\item عوارض \lr{RDS}: پنوموتوراکس، \lr{PDA}، برونکوپولمونری دیسپلازی (\lr{BPD}).
\end{itemize}
\subsection*{پیشگیری و درمان \lr{RDS}}
\begin{itemize}
\item تجویز گلوکوکورتیکوئید به مادر پیش از زایمان: \NRmark{بتامتازون}؛ در چند دوز، حداقل ۴۸ ساعت قبل از زایمان.
\item تجویز سورفاکتانت اندوتراکئال؛ ظرف چند ساعت اول تولد.
\item \lr{nasal CPAP}: تجویز زودهنگام.
\end{itemize}
\subsection*{شرایط درمان}
\begin{itemize}
\item \NRmath{PaO_2} بین \NRmark{۶۰–۷۰} میلی‌متر جیوه حفظ شود.
\item عبارت خط‌خورده: «\NRmath{SpO_2>90\%} حفظ شود».
\item اندیکاسیون ونتیلاسیون کمکی: شواهد نارسایی تنفسی؛
\NRmath{PCO_2>60}، \NRmath{pH<7.20}، \NRmath{PaO_2<50} با \NRmark{\NRmath{100\%\,O_2}}.
\item در \lr{RDS} آنتی‌بیوتیک وسیع‌الطیف برای ۴۸–۷۲ ساعت می‌دهیم، تا جواب \lr{BC} بیاید؛
\lr{R/O} سپسیس؛ \NRmark{آمپی‌سیلین ـ جنتامایسین}.
\end{itemize}

\NRpage{۶}
\subsection*{فلوچارت \lr{RDS}}
\begin{enumerate}
\item \lr{CPAP} اولیه.
\item تلاش تنفسی ادامه دارد، \NRmark{\NRmath{FiO_2<0.4}}: ادامهٔ \lr{CPAP} + \lr{ABG}.
\begin{itemize}
\item \NRmath{FiO_2<0.4}، \NRmath{pH\geq7.2}: ادامهٔ درمان.
\item \NRmath{FiO_2>0.4}، \NRmath{pH<7.2}: \NRmark{انتوبه} برای دادن \NRmark{سورفاکتانت} + \NRmark{\lr{ABG}}.
\end{itemize}
\item آپنه یا \NRmath{FiO_2\geq0.4}: انتوبه برای دادن سورفاکتانت + \lr{ABG}.
\item پس از سورفاکتانت:
\begin{itemize}
\item \NRmath{FiO_2\geq0.3}: دوباره سورفاکتانت.
\item \NRmath{FiO_2<0.3}، \NRmath{pH>7.25}: اکستوبه؛ \lr{CPAP} یا نازال \lr{IPPV}.
\end{itemize}
\end{enumerate}
\subsection*{کامپلیکیشن‌های \lr{RDS}: \lr{PDA}}
\begin{itemize}
\item در نوزادان نارس شایع است. در منبع، واژهٔ «خوش» خط خورده و «نارس» بالای آن نوشته شده است.
\item به طور طبیعی ۲۴–۴۸ ساعت پس از تولد بسته می‌شود.
\item درمان دارویی و جراحی دارد.
\item \lr{PDA} در فاز دیسترس: شانت شدید سیستمیک به پولمونری؛ منجر به \NRmark{\lr{HF}} می‌شود و \NRmark{ادم ریوی}.
\item نباید در این بیماران مایع‌درمانی شدید انجام بدهیم؛ در نهایت هایپرکاپنه، هایپوکسی و ادم ریوی.
\item \lr{PDA} ظرف روز ۲–۴ تولد علامت‌دار می‌شود؛ شانت چپ به راست؛ \NRmark{فشار نبض پهن می‌شود}.
\item نبض‌ها جهشی قابل لمس می‌شوند.
\item \lr{Murmur}: می‌تواند \NRmark{کانتینیوس} در \lr{S} و \lr{D}؛ یا کانتینیوس در \lr{S} \NRunclear{واژهٔ دوم نیز در منبع «کانتینیوس» خوانده می‌شود}.
\end{itemize}

\NRpage{۷}
\begin{itemize}
\item در \lr{PDA} به علت \lr{HF} و ادم ریه: \NRmark{رال + هپاتومگالی}؛ \NRmark{کاردیومگالی} در \lr{CXR}.
\item تشخیص \lr{PDA}: سونو دوبعدی و داپلر. مخفف در این سطرِ منبع به صورت \lr{RDA} نوشته شده است.
\end{itemize}
\subsection*{درمان \lr{PDA} طی \lr{RDS}}
\begin{itemize}
\item محدودیت حجم + دیورتیک.
\item اگر ظرف \lr{48 h} بیمار وارد بهبودی نشد: مهارکنندهٔ تولید \lr{PG}؛ ایبوپروفن، \NRmark{ایندومتاسین} یا استامینوفن.
\end{itemize}
\subsection*{کنترااندیکاسیون‌های ایندومتاسین}
\NRmark{ترومبوسایتوپنی زیر ۵۰٬۰۰۰، خونریزی، \NRmath{Cr>1.8}، الیگوری}.
\begin{itemize}
\item ۳۰\% نوزادان \lr{PDA} به دور اول درمان دارویی جواب نمی‌دهند.
\item ۲۰\% نوزادانی که دور اول جواب داده‌اند، دوباره عود می‌کنند.
\end{itemize}
\subsection*{عارضهٔ دوم \lr{RDS}: لیک هوا}
آمفیزم بینابینی ریه، پنومومدیاستن، آمفیزم زیرجلدی، پنوموتوراکس.
\par درمان لیک: سورفاکتانت و به حداقل رساندن لزوم تهویهٔ مثبت.
\subsection*{پنوموتوراکس در نوزادان}
\begin{itemize}
\item پس از احیای قلبی‌ریوی، پنومونی، مکونیوم آسپیریشن، هایپوپلازی ریوی، هرنی دیافراگماتیک.
\item \NRmark{پنوموتوراکس خودبه‌خودی با مالفورماسیون‌های کلیوی مرتبط است.}
\item در کمتر از ۱\% زایمان‌ها.
\end{itemize}

\NRpage{۸}
\subsection*{عارضهٔ سوم \lr{RDS}}
\NRmark{دیسپلازی برونکوپولمونری} (بیماری مزمن ریوی) (\lr{BPD}).
\begin{itemize}
\item \lr{BPD} یک تشخیص بالینی است.
\item \NRmark{وابستگی به اکسیژن} در ۳۶ هفتگی پس از لقاح، به همراه یافته‌های بالینی و تصویربرداری کاراکتریستیک؛ مربوط به ابنورمالیتی‌های آناتومیک.
\item اکسیژن بالای \NRmark{۴۰\%} برای ریهٔ نوزاد سمی است.
\item ونتیلاسیون \lr{high peak pressure} در نوزادان \NRunclear{واژهٔ بعدی شبیه «ایدو» است}: باروتروما.
\end{itemize}
\subsection*{سناریوی شایع \lr{BPD}}
\begin{itemize}
\item \NRmark{نوزاد دچار \lr{RDS} دچار عارضهٔ \lr{BPD} می‌شود.}
\item ممکن است با \NRmark{\lr{PDA}} و \NRmark{آمفیزم بینابینی} هم همراه شود.
\item کاراکتریستیک بیمار \lr{RDS} که دچار \lr{BPD} می‌شود:
\begin{itemize}
\item \lr{RDS} \NRmark{پس از ۲ هفته} رو به بهبودی نمی‌رود.
\item نیاز به ونتیلاسیون مکانیکی طولانی‌مدت.
\item نیاز به \NRmath{O_2} پس از هفتهٔ ۳۶ لقاح.
\end{itemize}
\item نوزادان کمتر از \lr{1000 gr} که نیاز به ونتیلاسیون دارند ولی \lr{RDS} ندارند: خطر بالای \lr{BPD}.
\end{itemize}
\subsection*{یافته‌های تصویربرداری در \lr{BPD}}
\begin{latin}
\begin{itemize}
\item \textbf{Lung opacification}.
\item Areas of overdistention.
\item Developing of \textbf{cysts}.
\item \textbf{Atelectasis} (spongelike appearance).
\end{itemize}
\end{latin}
\subsection*{پاتولوژی \lr{BPD}}
ادم بینابینی، متاپلازی مخاطی، فیبروز بینابینی، برونشیولیت نکروزان.

\NRpage{۹}
\subsection*{تظاهرات بالینی \lr{BPD}}
\begin{itemize}
\item وابستگی به \NRmath{O_2}.
\item \NRmark{هایپرکاپنه} + آلکالوز متابولیک جبرانی. در منبع «اسیدوز» خط خورده و «آلکالوز» بالای آن نوشته شده است.
\item فشار خون ریوی؛ \lr{poor growth}.
\item \NRmark{نارسایی قلب راست}؛ افزایش مقاومت راه هوایی.
\item در \lr{BPD} به صورت واکنشی انقباض برونش‌ها داریم؛ \NRmark{برونکودیلاتور} می‌دهیم.
\item در \lr{BPD}: رتراکشن شدید قفسهٔ سینه + \NRmark{کورپولمونل}؛ ریتنشن مایع؛ \NRmark{دیورتیک}.
\item موارد شدید \lr{BPD}: نیاز به ونتیلاسیون به مدت چندین ماه.
\end{itemize}
\subsection*{درمان \lr{BPD}}
\begin{itemize}
\item ونتیلاسیون: \NRmath{PaO_2\sim50\,\mathrm{mmHg}}؛ \NRmath{PaCO_2\sim50\text{--}75} (کمی بالاتر).
\item پرهیز حداکثری از باروتروما \NRunclear{واژهٔ میانی}.
\item \NRmark{دگزامتازون}: ریسک \lr{CP} و اختلالات نورودولوپمنتال.
\end{itemize}
\subsection*{عوارض طولانی‌مدت \lr{BPD}}
\begin{itemize}
\item تأخیر رشد؛ \lr{Reactive airways}.
\item نوزادان \lr{BPD} در ریسک ویروس \lr{RSV} یا ویروس سنسیشیال تنفسی؛ باید پروفیلاکسی بگیرند.
\end{itemize}
\subsection*{\lr{ROP}}
\begin{latin}
Retinopathy of prematurity (Retrolental fibroplasia).
\end{latin}
\begin{itemize}
\item سمیت \NRmath{O_2} باعث تشکیل عروق در شبکیه می‌شود.
\item در \NRmark{وزن تولد زیر \lr{1500 gr}} شایع‌تر است.
\end{itemize}

\NRpage{۱۰}
\begin{itemize}
\item هایپرکاپنه و هایپوکسی منجر به \lr{ROP} می‌شوند.
\item مراحل شدید \lr{ROP}: عروق در تمام رتین؛ ریسک جدا شدن شبکیه؛ \NRmark{لوکوکوریا} (به علت فیبروز پشت لنز)؛ \NRmark{گلوکوم}.
\item معمولاً \lr{ROP} \NRmark{دو چشم} را درگیر می‌کند و می‌تواند شدت علائم متفاوت باشد.
\item هیچ سطح \NRmath{O_2} برای \lr{ROP} ایمن نیست ولی \NRmath{O_2} شریانی \NRmark{۵۰–۷۰} باید باشد.
\item نوزاد نارس زیر ۳۲ هفته باید \NRmark{پس از ۴ هفته از تولد} یا پس از \NRmark{۳۴ هفتگی} از نظر \lr{ROP} غربالگری شوند \NRunclear{عدد دوم در منبع ۳۴ خوانده می‌شود؛ احتمال خوانش ۳۱ وجود دارد}.
\end{itemize}
\subsection*{\lr{TTN}}
\begin{itemize}
\item خودمحدودشونده است.
\item علائم: \NRmark{تاکی‌پنه}، رتراکشن خفیف، \NRmark{هایپوکسی}، گرانتینگ.
\item معمولاً \NRmark{علائم دیسترس شدید ندارند}.
\item ممکنه در \lr{TTN} \NRmark{سیانوز} دیده شود؛ تهویه با \NRmath{O_2} ۳۰–۴۰\%.
\item \lr{TTN} در نوزادان \NRmark{پره‌ترم} یا نوزادان با زایمان سخت یا با \lr{C/S} بدون لیبر.
\item \lr{TTN} در نوزادان با مادر \NRmark{دیابتی} یا مادر \NRmark{آسم آتوپیک} شایع‌تر است.
\item علت آن \NRmark{جذب کند مایع داخل ریه} است.
\end{itemize}

\NRpage{۱۱}
\subsection*{یافته‌های تصویربرداری در \lr{TTN}}
\begin{latin}
\begin{itemize}
\item Prominent central vascular markings.
\item \textbf{Fluid in the lung fissures}.
\item Small PE.
\end{itemize}
\end{latin}
محل: \NRmark{پری‌هیلر}.
\par عبارت خط‌خورده: \lr{Air bronchogram}؛ سطر خط‌خوردهٔ زیر آن: \NRunclear{متن زیر خط‌خوردگی قابل خواندن نیست؛ اصل آن در تصویر ۱۱ محفوظ است}.
\begin{itemize}
\item \NRmark{در \lr{TTN} الگوی رتیکولوگرانولر و \lr{air bronchogram} دیده نمی‌شود.}
\item اگر دیده شد: \lr{RDS} یا پنومونی است.
\end{itemize}
\subsection*{\lr{MAS} (سندرم مکونیوم آسپیریشن)}
\begin{itemize}
\item بیشتر در نوزادان ترم و \NRmark{پست‌ترم} دیده می‌شود.
\item وجود مکونیوم در مایع آمنیوتیک در جنین با \NRmark{پرزنتاسیون بریچ} شایع‌تر است.
\item وجود مکونیوم در مایع آمنیوتیک: دیسترس جنینی؛ آسفیکسی، هایپوکسی، \NRmark{اسیدوز}.
\item مکونیوم آسپیریشن ریسک \NRmark{پنومونی} و \NRmark{پنوموتوراکس} را بالا می‌برد.
\end{itemize}
\subsection*{علائم فوری مکونیوم آسپیریشن (\lr{MAS})}
مخففِ این عنوان در دست‌نویس به صورت \lr{MAP} آمده است.
\begin{itemize}
\item \NRmark{تاکی‌پنه}، هایپوکسی، هایپرکاپنه، انسداد مجاری هوایی کوچک؛ \lr{Causing ball-valve effect}.
\item \NRmark{پرهوايی}، لیک هوا، خارج آلوئول.
\item انسداد کامل مجاری هوای کوچک: \NRmark{آتلکتازی}.
\item ظرف \lr{24--48 h}: پنومونیت شیمیایی ایجاد می‌شود.
\item مکونیوم \NRmark{سورفاکتانت را غیرفعال می‌کند}.
\item \NRmark{فشار خون ریوی} در \lr{MAS} شایع است.
\end{itemize}

\NRpage{۱۲}
\begin{itemize}
\item افزایش مقاومت شریان‌های ریوی: ایجاد \NRmark{شانت راست به چپ}.
\end{itemize}
\subsection*{یافته‌های \lr{CXR}}
\NRmark{انفیلتراسیون پچی}، \NRmark{پرهوايی}، تخت شدن دیافراگم، افزایش قطر قدامی ـ خلفی، ریسک بالای پنوموتوراکس و پنومومدیاستن.
\subsection*{درمان}
\begin{itemize}
\item \NRmark{پیشگیری}، حمایت تنفسی.
\item اگر شواهد \lr{PPHN}: درمان خاص خود.
\item اگر هایپوکسی شدید به درمان با \NRmath{O_2} جواب نداد: درمان با سورفاکتانت، \lr{NO} استنشاقی، \lr{ECMO}.
\end{itemize}
\subsection*{\lr{PPHN}}
\begin{itemize}
\item در نوزادان \NRmark{پست‌ترم، ترم} و نزدیک به ترم.
\item \NRmark{هایپوکسی شدید} بدون شواهد بیماری پارانشیم ریه یا بیماری ساختاری قلب.
\item \lr{PPHN} معمولاً همراه با آسفیکسی یا \NRmark{مایع آمنیوتیک آغشته به مکونیوم} دیده می‌شود.
\item در \lr{PPHN}، اغلب \lr{CXR}: \NRmark{ریه نرمال است}.
\end{itemize}
\subsection*{علل ثانویهٔ \lr{PPHN}}
عفونت مادرزادی، پلی‌سیتمی (هایپرویسکوزیتی)، هایپوپلازی ریوی، هرنی دیافراگماتیک مادرزادی، بیماری \NRmark{سیانوتیک} مادرزادی قلب، هایپوگلایسمی، هایپوترمی.
\subsection*{شواهد کاراکتریستیک \lr{PPHN}}
\begin{itemize}
\item \NRmark{شانت راست به چپ شدید}: \lr{PDA}، فورامن اووال باز، اینتراپولمونری کانال‌ها.
\item هایپوکسی و اسیدوز رخ می‌دهد.
\end{itemize}

\NRpage{۱۳}
\begin{itemize}
\item \lr{PPHN} هایپوکسی ایجاد می‌کند؛ به علت هایپوکسی سایر بافت‌ها آسیب می‌بینند.
\item آسیب میوکارد: \lr{HF}، ناکفایتی گذرای میترال، انفارکت میوکارد.
\item عوارض \lr{PPHN}: \NRmark{ترومبوسایتوپنی، ترومبوز داخل قلب، آمبولی ریوی}.
\item تأیید تشخیص \lr{PPHN}: \NRmark{اکوکاردیوگرافی}.
\end{itemize}
\subsection*{درمان}
\begin{itemize}
\item حمایتی، اصلاح فشار خون پایین، آنمی و اسیدوز.
\item اگر اختلال میوکارد وجود دارد: دوپامین یا دوبوتامین.
\item \NRmark{مهم‌ترین درمان ونتیلاسیون کمکی است.}
\item \NRmark{سورفاکتانت در درمان \lr{PPHN} جایگاهی ندارد.}
\item اگر تنفس بهبود نیافت: \NRmark{\lr{NO} استنشاقی} + داروهای وازودیلاتور انتخابی شریان‌های ریوی.
\end{itemize}
\subsection*{\lr{Apnea of Prematurity}}
\begin{itemize}
\item آپنه یعنی قطع جریان هوا تنفسی برای \NRmark{بیش از ۱۰–۲۰ ثانیه}.
\item آپنهٔ طول‌کشیده معمولاً همراه با \NRmark{برادی‌کاردی} است.
\end{itemize}
\subsubsection*{انواع آپنه}
\begin{itemize}
\item مرکزی: حرکات قفسهٔ سینه وجود ندارد.
\item انسدادی: حرکات قفسهٔ سینه وجود دارد.
\item میکس: \NRmark{شایع‌ترین فرم}.
\end{itemize}
\begin{itemize}
\item آپنهٔ ایدیوپاتیک در نوزاد \NRmark{پره‌مچور}، در \NRmark{هفتهٔ اول}.
\item معمولاً تا هفتهٔ ۳۹–۴۰ پس از لقاح رفع می‌شود \NRunclear{عدد نخست در منبع ۳۹ خوانده می‌شود}.
\item \NRmark{نوزاد پره‌ترم به هایپوکسی به صورت پارادوکسیکال جواب می‌دهد.}
\item در نوزاد پره‌ترم ضعف عضلات حنجره وجود دارد؛ انسداد راه هوایی.
\end{itemize}

\NRpage{۱۴}
درمان اصلی آپنه: \NRmark{متیل‌گزانتین‌ها} (کافئین، تئوفیلین).

\clearpage
\section*{پیوست: تصاویر کامل منبع}
تمام متن اصلی، حاشیه‌ها، خط‌خوردگی‌ها، فلش‌ها، هایلایت‌ها و سربرگ‌ها در تصاویر زیر محفوظ‌اند.
\IfFileExists{images/pic01.jpg}{\def\NRimagebase{images/}}{%
  \def\NRimagebase{pediatrics/ped-neonatology/respiratory diseases of neonates/images/}}
\newcommand{\NRoriginal}[2]{%
  \clearpage\subsection*{اصل تصویر #1: \lr{\texttt{pic#2.jpg}}}%
  \begin{center}%
  \includegraphics[width=\linewidth,height=0.87\textheight,keepaspectratio]{\NRimagebase pic#2.jpg}%
  \end{center}}
\NRoriginal{۱}{01}
\NRoriginal{۲}{02}
\NRoriginal{۳}{03}
\NRoriginal{۴}{04}
\NRoriginal{۵}{05}
\NRoriginal{۶}{06}
\NRoriginal{۷}{07}
\NRoriginal{۸}{08}
\NRoriginal{۹}{09}
\NRoriginal{۱۰}{10}
\NRoriginal{۱۱}{11}
\NRoriginal{۱۲}{12}
\NRoriginal{۱۳}{13}
\NRoriginal{۱۴}{14}
\NRfinish
